\section{Cut, parity and modular templates}\label{sec:templates}

Every obstruction in this and the next sections concerns an additional
restriction on a colouring or on a proposed algorithm. We supply
actual legal colourings or cite an already established family result
where needed to keep that distinction explicit.

\subsection{A maximum-cut template with no exceptions}

The template chooses a maximum cut $(S,T)$, requires all edges in
$G[S]$ to have weight zero, and requires $R(v)\equiv\mathbb{1}_T(v)$
at every vertex. For any vertex set $U$ the actual handshake identity is
\begin{equation}\label{eq:boundaryhandshake}
 \sum_{v\in U}R(v)=2\sum_{e\in E(U)}x_e+
 \sum_{e\in\delta(U)}x_e.
\end{equation}
With parity-exception set $X$, the last sum has parity
$|U\cap T|+|U\cap X|$. A single bridge of odd required parity
must therefore be mixed. This is a necessary boundary constraint,
not a guarantee of an extendible interface.

\begin{proposition}\label{prop:cutobstruction}
For every odd $r\ge3$, take a triangle $abc$ and attach one copy
of $K_{r,r}$ to each triangle vertex by a bridge to a distinguished
module root. The resulting graph $G_r$ has minimum degree three
and a legal doubled two-colouring, but no maximum cut admits the
zero-exception template above. Exactly one exception suffices.
\end{proposition}
\begin{proof}
The graph has $3+6r$ vertices and $3r^2+6$ edges. A maximum cut
has size $3r^2+5$, since each module contributes at most $r^2$,
each bridge at most one, and the triangle at most two; orienting
the module bipartitions independently attains all these bounds.
Every maximum cut therefore cuts every module edge and every bridge.
Each module has $r$ vertices in $T$, odd. With no exceptions,
\eqref{eq:boundaryhandshake} forces each module bridge to have
weight 1. The total red sum is even; the three modules contribute
an odd number of $T$ vertices, so precisely one triangle vertex is
in $T$. The other two triangle vertices lie in $S$, their common
edge has weight zero, and their edges to the $T$ vertex must both
have weight 1 to give them even red sums after the mixed bridges.
They both have counts $(2,4)$, contradicting their common blue edge.

Here is a full legal colouring. Give $ab,bc,ca$ weights $0,1,2$
and all bridges weight 1. Then
$a:(3,3), b:(2,4), c:(4,2)$. In each module use rows and columns
$0,\ldots,r-1$, with root at row zero. In the $a$ module give the
root row all weights 1 and all other entries zero; the root is
$(r+1,r+1)$, other rows $(0,2r)$ and columns $(1,2r-1)$.
For the $b$ module start with root entry $(0,0)$ of weight 1,
row 1 entries in all other columns of weight 1, all remaining
entries zero, and complement its internal weights. The root is
then $(2r,2)$; row counts are even and column counts odd. For the
$c$ module give the root row zero, every other row weight 1 only
in column zero, and all other entries zero. Its root is
$(1,2r+1)$; other rows $(1,2r-1)$, column zero $(r-1,r+1)$,
other columns $(0,2r)$. Row--column parity verifies every module
edge and the displayed root pairs verify the bridges. The triangle
edges are checked directly.

Choose the maximum cut with $a,b\in T,c\in S$, module roots on
opposite sides to their triangle neighbours. The module patterns
have precisely the cut parities, as do $a,c$. Only $b$ is an
exception. Thus the minimum exception count in this template is one.
\end{proof}

The same graph also excludes the different prescription
$R(v)\equiv d_G(v)\pmod2$ at \emph{every} vertex, without mentioning
a cut. Each module has odd sum of original degrees, so its bridge
would again be mixed. The three triangle vertices would then have
distinct odd red degrees, necessarily $1,3,5$, because they have
equal original degree three. Subtracting the mixed bridge contribution
would require internal triangle red sums $0,2,4$, impossible at the
edge between the zero and four vertices. This new restriction also
fails despite the explicit legal colouring in the proof.

\subsection{Small boundary modules and parity transfer}

There is a seven-vertex module consisting of triangle $abc$,
two-edge pendant paths at $b,c$, and a fixed mixed outside edge
at $a$. If its internal cut is maximum, all four path edges cross
and two triangle edges cross. Under the zero-exception template,
the first edge of a two-edge pendant must be mixed: its leaf and
triangle vertex lie on the same side, and the middle parity forces
that first weight odd. The module has an odd number of $T$ vertices
by its mixed boundary, hence one triangle vertex in $T$. Each of
the two $S$ vertices has one mixed outside contribution, so its
edge to the $T$ vertex is mixed too. Their common blue edge again
joins counts $(2,4)$. Thus fixed mixed boundaries need not extend
even when total parity is compatible.

Put two-edge pendant paths at all three triangle vertices to obtain
a nine-vertex graph $O$. Every maximum cut and zero-exception
template fails by the same argument. For clarity, in edge order
$ab,ac,bc,aa_1,a_1a_2,bb_1,b_1b_2,cc_1,c_1c_2$ a legal word is
$(1,2,0,1,0,0,0,1,1)$. With
$T=\{b,c,a_1,b_2,c_2\}$ it has exactly the one exception $b_2$.
This small existence statement is subsumed by the stronger
minimum-degree-three family above, but its boundary properties differ.

In $O$ choose root $a_2$ and give it arbitrary outside total red
contribution $A$. Internal maximum-cut template feasibility with no
exceptions requires $A$ odd. If the pendant at $a$ instead has
length four and its endpoint is the root, feasibility requires $A$
even. To see both assertions, for a pendant of length $2t$ apply
\eqref{eq:boundaryhandshake} to its $2t$ nontriangle vertices.
Its first weight $q$ satisfies $q+A\equiv t\pmod2$. The wrong
boundary parity in either case forces $q=1$ and one triangle vertex
in $T$, repeating the $(2,4)$ conflict. These are necessary
conditions only; sufficiency of every boundary with that parity is
not claimed.

\subsection{A six-vertex modulo-three obstruction}

Take triangles $uab$ and $vcd$ joined by bridge $uv$. Even though
this cactus is three-colourable, one cannot require every present
red and blue comparison to be nonzero modulo three. Write the left
weights $ua,ub,ab,uv$ as $p,q,t,z$ and put $h=z-t\pmod3$.
The degree-two edge $ab$ forces $p\ne q$. On $ua$ the restrictions
are respectively $q+h\ne2$, $q+h=1$, $q+h\ne0$ when
$p=0,1,2$. The $ub$ conditions interchange $p,q$. Unordered pairs
$\{0,1\}$ and $\{1,2\}$ immediately contradict these restrictions;
only $\{p,q\}=\{0,2\}$ remains. The right triangle obeys the
same rule, so $R(u)\equiv R(v)\equiv2+z$ and
$B(u)\equiv B(v)\equiv1-z$. Every bridge weight then fails.
In contrast the actual legal word in order
$ua,ub,uv,ab,vc,vd,cd$ is $(1,2,2,0,2,0,1)$, giving counts
$u:(5,1),a:(1,3),b:(2,2),v:(4,2),c:(3,1),d:(1,3)$.
For instance a legitimate red degree difference of three need not
survive modulo three. The modulo-two certificate has its own
independent restrictions in Section~\ref{sec:mod2limits}.

\section{Barriers to nonincreasing local descent}\label{sec:descent}

\subsection{A closed two-state single-edge plateau}

On the wheel with centre 0 and rim $123451$, use edge order
\[
 01,02,03,04,05,12,15,23,34,45.
\]
The words
$x=(1,2,1,1,0,0,2,2,0,0)$ and
$y=(1,2,1,1,0,0,2,2,2,0)$ both have $\Psi=1$.
All single-edge changes, allowing even direct $0\leftrightarrow2$
changes, have the following conflict values. Each triple gives the
result for new weights $0,1,2$.
\begin{center}\small
\begin{tabular}{ccc}\toprule Edge&From $x$&From $y$\\\midrule
$01$&$(4,1,4)$&$(4,1,4)$\\
$02$&$(2,4,1)$&$(2,2,1)$\\
$03$&$(2,1,4)$&$(4,1,4)$\\
$04$&$(4,1,4)$&$(4,1,2)$\\
$05$&$(1,2,2)$&$(1,4,2)$\\
$12$&$(1,3,2)$&$(1,5,2)$\\
$15$&$(2,5,1)$&$(2,3,1)$\\
$23$&$(4,3,1)$&$(2,2,1)$\\
$34$&$(1,4,1)$&$(1,4,1)$\\
$45$&$(1,2,2)$&$(1,3,4)$\\\bottomrule
\end{tabular}
\end{center}
The only nontrivial nonincreasing steps exchange $x,y$. By induction
every such process remains in this two-state set. A simultaneous
change of $01$ from 1 to zero and $12$ from zero to 1 gives
$(0,2,1,1,0,1,2,2,0,0)$ with $\Psi=0$. The complete finite table
proves a barrier for single-edge monotonicity, including neutral
steps; it does not exclude larger simultaneous repairs.

\subsection{Strict positive minima after fixed-target optimisation}

On a $d$-regular graph every assignment realising $q$ has
\[
 F(q)=2|\{uv:q(u)=q(v)\}|.
\]
Suppose equal-target edges are a perfect matching, every available
target $q(v)\pm1$ occurs on at least three neighbours, and every
available $q(v)\pm2$ on at least two. Then $F(q)=n>0$ and
every realisable target at $\ell_1$ distance two has larger $F$.
Changing one coordinate by two creates at least two equalities and
removes at most one. Changing two coordinates by one creates at
least two equalities to unchanged neighbours at each coordinate,
removing at most the two old matching edges. Thus no neutral first
step is available, even after completely optimising each fibre.

Here is an explicit infinite realisable family for that criterion.
For any even $N\ge4$, take vertices $(i,t)$ with $0\le i\le26$
and $t\in\mathbb Z_N$, target $q(i,t)=i$. Put a perfect matching
within each layer by pairing consecutive even/odd $t$ labels. Between
layers start with $a_{ij}=3$ when $|i-j|=1$, 2 when $|i-j|=2$,
and zero otherwise. Add each increment in the first table below and
its reflection $(26-j,26-i)$, then make $a$ symmetric.
\begin{center}\small
\begin{tabular}{rrr@{\qquad}rrr}\toprule
$i$&$j$&increment&$i$&$j$&increment\\\midrule
0&1&1&0&2&1\\0&4&2&0&11&2\\0&24&1&1&3&2\\
1&20&1&5&14&2\\6&13&1&7&9&2\\8&10&2&&&\\\bottomrule
\end{tabular}
\end{center}
Its rows sum to 12 and $0\le a_{ij}\le4$. For $i<j$, add the
$a_{ij}$ matchings $(i,t)(j,t+s)$, $s=0,\ldots,a_{ij}-1$.
They are distinct for $N\ge4$. Together with the within-layer edge
the graph is 13-regular. Shifts zero and one between consecutive
layers make it connected. Equal targets occur exactly on the
within-layer perfect matching and the required neighbour counts hold.

To realise $q$, for $i<j,i+j\le26$ set $b_{ij}=0$ except
\[
\begin{array}{c|rrrrrrrrr}
(i,j)&(3,4)&(5,7)&(6,7)&(8,9)&(9,11)&(10,11)&(10,12)&(11,13)&(12,14)\\
b_{ij}&2&3&4&6&1&6&4&4&2
\end{array}
\]
and define $b_{ij}=2a_{ij}-b_{26-j,26-i}$ for $i+j>26$;
symmetrise. Within-layer weights at $i=0,\ldots,13$ are
\[
 (c_0,\ldots,c_{13})=(0,1,0,1,2,2,0,0,2,2,0,0,2,1),
\]
with $c_{26-i}=2-c_i$. Direct row addition verifies
$0\le b_{ij}\le2a_{ij}$ and $c_i+\sum_{j\ne i}b_{ij}=i$.
Give shift-$s$ edges between layers weight
$\min\{2,\max\{0,b_{ij}-2s\}\}$ and within-layer edges weight
$c_i$. Each vertex receives total $i$, as required. The tables are
a fixed integer certificate, independently checked in the artifact.
Thus for this connected infinite family of order $27N$,
$F(q)=27N$ and every two-unit neighbour has $F\ge27N+2$.
These regular graphs have other legal targets by~\cite{GOW26}.
This bad initial target cannot itself arise by subtracting degrees
from a proper 1--2--3 neighbour-sum weighting, since its equal-target
matching edges have equal original degrees.

\subsection{A proper 1--2--3 initialisation can also be trapped}

On the five-vertex path use $x=(1,0,0,1)$. Then $w=x+1$ has
proper neighbouring sums $(2,3,2,3,2)$, but the doubled target
$q=(1,1,0,1,1)$ has $F(q)=2$. A path target determines its weights
uniquely by leaf recursion. Its alternating bipartition sum is zero,
so a two-unit target change must be two opposite-balanced unit
changes, corresponding to an alternating contiguous path of weight
changes. The two adjacent internal zero weights preclude all such
paths except the following eight:
\begin{center}\small
\begin{tabular}{cc@{\qquad}cc}\toprule
New word&$F$&New word&$F$\\\midrule
$(0,0,0,1)$&3&$(2,0,0,1)$&3\\
$(1,1,0,1)$&3&$(1,0,1,1)$&3\\
$(1,0,0,0)$&3&$(1,0,0,2)$&3\\
$(0,1,0,1)$&5&$(1,0,1,0)$&5\\\bottomrule
\end{tabular}
\end{center}
No nonincreasing two-unit first step exists, even with an initially
proper 1--2--3 weighting. The word $(0,0,1,1)$ is legal and has
target distance four from $q$, attaining the least possible
nonincreasing first-step distance. This does not prove that distance
four is sufficient for general graphs or specially selected initialisations.

\section{Fixed-target and radius obstructions}\label{sec:targetbarriers}

\subsection{Bridge locking, capacity and pure odd defects}

For a bridge $e=uv$ whose $u$ side $U$ is bipartite with signs
$\varepsilon$, fixed internal target demands give
\[
 x_e=\varepsilon(u)\sum_{v\in U}\varepsilon(v)q(v).
\]
Internal contributions cancel. Thus a bridge may be fixed uniquely;
on a forest the entire weighting is unique for a fixed target.
Take a triangle $uab$ with path $u w l$, in edge order
$ua,ub,uw,ab,wl$. The word $(0,1,1,1,1)$ is realisable and
comes from a proper 1--2--3 weighting, but the bridge $uw$ joins
equal red degree two, so legality would require $x_{uw}=0$.
Its bipartite side forces $x_{uw}=q(w)-q(l)=1$. No fixed-target
repair exists. The same graph has the legal word $(1,0,0,2,0)$.

For a purely integer obstruction, take triangles $012,345$, bridge
$03$, and leaves $26,57,38$. In edge order
\[
 01,02,12,34,35,45,03,26,57,38
\]
use $(0,1,1,0,1,1,1,2,2,0)$. Its targets and blue degrees are
\[
 q=(2,1,4,2,1,4,2,2,0),\quad B=(4,3,2,6,3,2,0,0,2).
\]
The sums $q+d=(5,3,7,6,3,7,3,3,1)$ distinguish all neighbours.
Only bridge $03$ is forced blue, and there are no forced red edges.
After deleting that bridge, the two free components each have target
sum nine. At $S=T=\varnothing$ the factor certificate is
$g=0,o=2$. However giving each triangle the word
$(1/2,3/2,1/2)$, setting the bridge zero and retaining the leaf
weights is a continuous legal realisation of exactly $q$. Hence
continuous feasibility at integer targets need not permit integer
rounding. A legal integer repair changes $02,35$ to 2 and the
bridge to zero, changing the targets at 2 and 5 by one. The actual
new red degrees are $(2,1,5,2,1,5,2,2,0)$.

The parity lower bound alone can also be strict. On edge set
$01,02,03,12,34$ the word $(1,0,0,0,0)$ has all free-component
target sums even, but vertex 1 needs red degree one through its
only free neighbour 2 of red target zero. The capacity pair
$S=\{2\},T=\{1\}$ has $g=-1$. Thus its quotient join bound is
zero while $F(q)=1$.

\subsection{Failure of exchange convexity}

For a triangle the targets uniquely determine weights by
$x_{ij}=(q_i+q_j-q_k)/2$. The targets $(1,2,3)$ and $(3,2,1)$
are legal. Starting from the first toward the second with step
$e_0$, the possible second steps are $e_0$ and $-e_2$.
The first gives target pair $(3,2,3),(1,2,1)$, each with $F=2$;
the second gives two copies of $(2,2,2)$, each with $F=6$.
Neither satisfies the fixed-cost exchange inequality with initial
total zero. In particular the midpoint of two legal targets can
be realisable with positive conflict. The favourable classical
exchange properties in Section~\ref{sec:targets} cannot be assigned
to the changing conflict-cost function $F$.

\subsection{An unbounded fixed-parity repair radius}\label{sec:radius}

\begin{proposition}\label{prop:radius}
Let even $c\ge2$ and $k\ge1$. Take $c$ disjoint cycles of length
$L=4k+2$ and two nonadjacent hubs $a,b$, each joined to every cycle
vertex. Put $N=cL$. Give the hubs target zero and every cycle
vertex target four, realised by cycle weights 2 and hub weights
zero. Fix the maximum cut with both hubs in $S$ and alternate
$S,T$ around each cycle, and require final red parity zero on $S$,
one on $T$. The least maximum coordinate change from those targets is
\[
 \rho_{\min}=2\lceil c/4\rceil,
\]
and this bound is achieved by legal doubled colourings. With $E$
parity exceptions and maximum change at most $M$, necessarily
$2M+E\ge c$.
\end{proposition}
\begin{proof}
The cut has size $2N$ and is maximum. If hubs are on opposite sides,
the hub edges contribute $N$ and cycle edges at most $N$. If hubs
share a side and $t$ cycle vertices share it, hub edges contribute
$2(N-t)$ and cycle edges at most $2t$. The alternating cut reaches
$2N$. Each cycle has $2k+1$ odd-labelled vertices. Its red boundary
sum is therefore a positive odd integer by
\eqref{eq:boundaryhandshake}. Hence $R(a)+R(b)\ge c$. Both hub
targets have required even parity, forcing the larger to be at
least $2\lceil c/4\rceil$.

Write each cycle $s_0,t_0,\ldots,s_{2k},t_{2k},s_0$, starting
with all cycle weights 2. Change one hub edge to $t_0$ to weight
1. For $i=1,\ldots,k$ change both edges incident with $s_{2i}$
along its two adjacent $t$ vertices to 1. The counts are
$t_0:(5,3)$, other $t$ vertices $(3,5)$, even positive-index
$s$ vertices $(2,6)$, and other $s$ vertices $(4,4)$.
Assign $2\lfloor c/4\rfloor$ mixed hub edges to $a$ and the rest
to $b$. Hub red degrees are those two even numbers, with blue
degrees at least $2N-c\ge22$. All cycle and mixed hub comparisons
are parity-separated. The pure blue hub edges to same-side $s$
vertices compare a degree at least 22 to one at most six. The
maximum change is exactly the asserted radius.

With $E$ exceptions, at least $c-E$ cycles have no exception and
each still needs at least one boundary red copy. The two hub red
degrees are each at most $M$, giving $2M\ge c-E$ independently
of any local irregularity requirement.
\end{proof}

This is a minimum-degree-four, 2-connected family with no peripheral
trees or degree-two chains. The starting target is feasible but is
not a legal colouring. The radius conclusion fixes the stated cut
parities. Without fixing them, keep hub weights zero and use the
cyclic word $(0,1,2,0,1,2)(0,1,1,0)^{k-1}$ on each cycle.
Positions two apart differ, so these words are legal; all changes
from the initial target are at most four, with unrestricted parity
exceptions. Thus the lower bound is not a distance statement about
all legal colourings.

Even an everywhere-one centre window is too narrow on $C_5$:
some red degree in any legal colouring is zero or four. If all red
degrees were in $\{1,2,3\}$, the number of degree-two vertices
would be odd by the handshake sum, and at most two by independence,
hence exactly one. The other four form an alternating 1--3 path.
An internal degree-one vertex between two degree-three vertices
would need two incident weights each at least one, impossible.
The word $(0,0,1,1,2)$ supplies a legal colouring. Translating all
degrees by fixed equal external offsets gives the corresponding
fixed-interface small-window obstruction.

\section{Sharp interface limits and multiple marks}\label{sec:interfacebarriers}

The path-length-six threshold in Theorem~\ref{thm:decorated} is
sharp for arbitrary fixed physical endpoint data. On a five-edge
bare path with end weights 1 and both endpoint counts $(3,3)$,
write the word $(1,a,b,c,1)$. Each mixed end edge forces $a=c=1$.
The middle edge then joins identical counts $(1+b,3-b)$ for any
$b$, and fails. The data are physically realised by adding two
mixed leaf edges at each endpoint. That whole ten-vertex tree is
colourable by Lemma~\ref{lem:treeports}; only the fixed interface
is impossible. The tree and clique capacity obstructions of
Sections~\ref{sec:treecapacity} and~\ref{sec:cliquecapacity} independently
show sharpness of their preserved-colour thresholds.

An arbitrary-root-sum interface also fails for general graphs. On
triangle $vabv$ with leaf $vz$, requiring the internal root red
sum zero forces all three incident weights zero. The endpoints of
$ab$ then have identical red and blue degrees whatever its weight.
Sum six is impossible by complementation. Nonetheless the word
$(x_{va},x_{vb},x_{ab},x_{vz})=(0,2,1,0)$ is legal, with
$(R(v),R(a),R(b),R(z))=(2,1,3,0)$ and blue degrees $(4,3,1,2)$.
The full clique interface is stronger than ordinary colourability.

Multiple fixed bridges need not preserve the one-bridge candidate
counts. On rooted path $r-u-v$, give $r$ external counts $(2,2)$
from two mixed bridges. If the internal word is $(a,b)$, the red
degrees are $(2+a,a+b,b)$ and blue degrees $(4-a,4-a-b,2-b)$.
At $a=0$ the edge $uv$ forces $b=0$, then $ru$ fails in blue;
at $a=2$ it forces $b=2$, then $ru$ fails in red. At $a=1$
the first edge excludes $b=0,2$, leaving exactly $(1,1)$.
Thus the internal root candidate set is the singleton $\{1\}$.

Finally consider a cubic-root safe-core requirement that every root
red degree belong to $\safe=\{0,1,3,5,6\}$ and every marked
mixed port have root red degree in $\{1,5\}$. It cannot accept
arbitrary many marks. For a two-vertex core with one mixed mark
and one pure port at each vertex, a central edge of weight 0 forces
both red sums 1, weight 1 allows neither strong marked sum, and
weight 2 forces both red sums 5. No choice is legal.

The failure persists with four independent marked roots: take a
double-star core with centres $u,v$, two leaves at each centre,
and one mixed and one pure port at every leaf. A strong marked leaf
forces its core edge and its pure port to be both zero (red sum 1)
or both 2 (red sum 5). If the central edge has weight zero, safety
forces both centre red sums zero; if it has weight 2, safety forces
both centre sums six. Both cases conflict on the central edge.
If its weight is 1, each centre with two zero leaf edges has sum
1, conflicting with those leaves, and with two 2-edges has sum 5,
also conflicting. Its only possible choice is one of each, giving
sum three at both centres, again a central conflict. This is a
complete parameter argument; its finite assignment domain has
$3^5\cdot2^4=3888$ members. The weak two-mark lemma in A1 imposes
different hypotheses and permits certain unexposed roots to have
red sums 2 or 4. These strong-mark obstructions therefore do not
contradict the subcubic proof.
